% !TEX root = ../main.tex
\chapter{Several particles (Mehrere Teilchen)}\label{ch:c15}
\begin{paracol}{2}

\begin{sources}
\begin{tabularx}{\linewidth}{@{}l l L@{}}
\toprule
\textbf{Video} & \textbf{Length} & \textbf{Content}\\
\midrule
\seg{15}{1}{Lecture 15: Several particles} & 1:51:40 & Recall of the one-particle Lagrangian; the configuration space of $N$ particles as the Cartesian product of $N$ copies of absolute space and its smooth atlas; the motion of $N$ particles as one curve in the configuration space (and back, by projections); composite systems without interaction: the Lagrangian as the sum of the one-particle Lagrangians on the tangent bundle of the configuration space; an example with two potentials; interaction between systems as a further term; a worked example: two free particles, with the Euler--Lagrange equations computed in the product chart. The recorded date is 12 June 2014.\\
\bottomrule
\end{tabularx}

\smallskip
\textbf{Transcript} \file{materials/transcripts/C15-V01.txt} (German YouTube auto-captions, mediocre in the spoken formulas; the formulas are reconstructed by recomputing and marked [?] where unsure).
\end{sources}

\begin{reading}[Where this lecture sits in the other notes]
Not treated in the other courses.
Continues \cref{ch:c14} (Euler--Lagrange equations, Lagrangian of one particle); the next lecture takes up the interactions and the constraints.
\end{reading}

% ==================================================================
\section{Who lectures, and the plan}\label{sec:mp-plan}

\begin{coursenote}[A substitute lecturer (as said in class)]\label{cn:mp-benni}
The lecture is not given by Schuller (the header of the transcript says so, but the speaker introduces himself): ``many of you may know me from the tutorial, I am Benni, I stand in for Frederik, who is travelling and does not want to let any lecture fall out, and does not want to impose on you to come in once more on a Saturday'', so he gives it today; ``you will have to put up with me''. He asks whether he is understood without a microphone, and whether his handwriting can be read (``I am no calligrapher; if you cannot read it just ask me what it means'') \ts{15}{1}{0:00}--\ts{15}{1}{1:58}. In these notes ``the lecturer'' is this substitute. He also mentions several times that the missing proofs are for Frederik to explain \emph{next lecture} (\cref{cn:mp-proof}).
\end{coursenote}

Today's topic is the Lagrangian of \emph{several} particles; above all the question how to assemble the Lagrangian of many particles from the one-particle Lagrangian of the previous lecture \ts{15}{1}{0:42}--\ts{15}{1}{1:27}. We begin with a recall.

\begin{recall}[The Lagrangian of one particle (\cref{ch:c14})]\label{rec:mp-one}
We looked at a particle of mass $m$ with a charge $c$; its \emph{world path} in absolute space, the (projected) curve $\gamma\colon\R\to\mathcal R$ into the absolute space $\mathcal R$, which carries a topology, an atlas and a metric $g$. The particle follows the \emph{stationary curve} of the action $S[\gamma]=\int\mathrm d\lambda\,L(\tilde\gamma(\lambda))$, the integral of the Lagrangian along the lifted curve. The Lagrangian is the function on the tangent bundle (the bundle projection $\pi$ sends the vector to its foot point)
\[
  L(X)\;=\;\tfrac12\,m\,g_{\pi(X)}(X,X)\;-\;c\,\Phi\bigl(\pi(X)\bigr),
\]
that is, one evaluates the metric at the foot point, inserts the vector twice, and subtracts the potential $\Phi$ of the interaction field, weighted by the charge $c$ and evaluated at the foot point, where the interaction field has to be \emph{conservative} (\cref{eq:var-Lcoupled}, \cref{def:nd-conservative}). (Capital $X$ is written by the lecturer with a ``hook'' so as to distinguish it from the small $x$.) \ts{15}{1}{2:18}--\ts{15}{1}{7:00}.
\end{recall}

The aim today is to answer two questions \ts{15}{1}{7:00}--\ts{15}{1}{8:56}:
\begin{enumerate}
  \item \emph{How do I construct the Lagrangian of $N$ particles from the Lagrangian of one particle?}
  \item \emph{How do I implement interactions between the different particles or systems?}
\end{enumerate}
That is the rough plan. Questions so far? A remark from the audience: ``we should name the parameters the same'' (``that is right''); the lecturer agrees \ts{15}{1}{9:22}--\ts{15}{1}{10:08} [?: the remark is hardly audible in the captions; probably about using the same curve parameter $\lambda$ for all particles, as is done below].

% ==================================================================
\section{The configuration space of $N$ particles}\label{sec:mp-config}

The first chapter: how do I build the \textbf{configuration space (Konfigurationsraum)} for $N$ particles out of the absolute space $\mathcal R$ on which the particles buzz around? So far we looked at one particle on this one space, and there the configuration space is the absolute space itself; but how do I build the space of all possible points that \emph{all} particles can take simultaneously? \ts{15}{1}{10:08}--\ts{15}{1}{11:17}. ``Has anyone an idea? It is actually quite simple'' \ts{15}{1}{11:40}.

\begin{definition}[New concepts: the configuration space of $N$ particles]\label{def:mp-config}
Let $\mathcal R$ be absolute space (a $d$-dimensional smooth manifold, $d=3$ in the application). The \textbf{configuration space (Konfigurationsraum)} of $N$ particles is the Cartesian product of $N$ copies of $\mathcal R$,
\begin{equation}\label{eq:mp-config}
  \mathcal Q\;:=\;\underbrace{\mathcal R\times\dots\times\mathcal R}_{N\text{ factors}}\;=\;\bigl\{(p_1,\dots,p_N)\ \big|\ p_j\in\mathcal R\bigr\},
\end{equation}
the set of $N$-tuples of points of $\mathcal R$ (``each of these $p$'s lies in this absolute space''), the $j$-th entry being the position of the $j$-th particle \ts{15}{1}{12:03}--\ts{15}{1}{13:55}.
\end{definition}
(The lecturer writes it with a capital curly letter, the symbol is garbled in the captions; $\mathcal Q$ here [?].)

As it stands it is only a \emph{set}. If we want to do mechanics on it, we need more structure: we would like $\mathcal Q$ to be a \emph{smooth manifold}, just as $\mathcal R$ is. How does one get a smooth manifold from a set? One needs a topology and a smooth atlas, that is, all the charts (``exactly: we must find the smooth atlas'') \ts{15}{1}{14:17}--\ts{15}{1}{15:27}. We want to \emph{inherit} the structure we have on $\mathcal R$, where we have charts, and let it pull up to $\mathcal Q$ \ts{15}{1}{15:48}.

Take charts $(U_1,x_1),\dots,(U_N,x_N)$ of the atlas of $\mathcal R$ (they need not be the same chart, ``these are all from the atlas of $\mathcal R$''). Define a map from the Cartesian product of the chart domains into $\R^{Nd}$ (``the dimension of the configuration space is $N$ times $d$'' \ts{15}{1}{16:08}--\ts{15}{1}{16:49}) by taking, as entries, the chart components of the successive points:

\begin{definition}[New concepts: the product chart $\omega$ on $\mathcal Q$]\label{def:mp-omega}
For charts $(U_j,x_j)$, $j=1,\dots,N$, of the atlas of $\mathcal R$ define
\begin{equation}\label{eq:mp-omega}
  \begin{aligned}
  &\omega\colon U_1\times\dots\times U_N\to\R^{Nd},\\
  &\omega(p_1,\dots,p_N)\;:=\;\bigl(x_1^1(p_1),\dots,x_1^d(p_1),\ x_2^1(p_2),\dots,x_2^d(p_2),\ \dots,\ x_N^1(p_N),\dots,x_N^d(p_N)\bigr).
  \end{aligned}
\end{equation}
That is, the components of $\omega$ are $\omega^1,\dots,\omega^{Nd}$ with
\begin{equation}\label{eq:mp-omegacomp}
  \omega^{\,d(j-1)+m}(p_1,\dots,p_N)\;=\;x_j^m(p_j),\qquad m=1,\dots,d,\ \ j=1,\dots,N
\end{equation}
(for $j=1$ the entries $1,\dots,d$, then for $j=2$ the entries $d+1,\dots,2d$, and so on up to the $N$-th: one takes the first component of the first chart on $p_1$, \dots, the $d$-th component of the first chart on $p_1$, then the first of the second chart on $p_2$, and so on) \ts{15}{1}{16:08}--\ts{15}{1}{25:14}.
\end{definition}
(``The chart $\omega$ has a bad name for the entries; here it is the symbol with which we denote the components''; the lecturer first calls the chart $\Omega$ and the components $\omega^1,\dots$, and mixes the two. Here the chart is $\omega$ and its components $\omega^k$.)

\begin{aside}[A joke (as said in class)]\label{as:mp-joke}
A remark by the lecturer while writing: ``I do it like in exams, copying from the neighbour'': no, ``what does this do? please do not'', officially marked as a joke, ``the tutor said I should copy it'' [?: the captions are garbled]. \ts{15}{1}{22:28}
\end{aside}

\begin{theorem}[The configuration space is a smooth manifold]\label{thm:mp-smooth}
The set $\mathcal Q$ with the product topology and the atlas $\mathcal A_{\mathcal Q}$ of all charts $(U_1\times\dots\times U_N,\omega)$ constructed from all $N$-tuples of charts of $\mathcal R$ is a smooth manifold of dimension $Nd$.
\end{theorem}
\begin{proof}
Topology: each factor $\mathcal R$ carries a topology, and we multiply them as Cartesian product: ``what we can do is use the product topology'' (as a set of open sets of the product) \ts{15}{1}{26:40}--\ts{15}{1}{27:27}. Atlas: the chart domains cover $\mathcal Q$, and each $\omega$ is a homeomorphism onto an open subset of $\R^{Nd}$ since each $x_j$ is one. The chart transition maps are, in components,
\[
  \omega'\circ\omega^{-1}\;=\;\bigl(x_1'\circ x_1^{-1}\bigr)\times\dots\times\bigl(x_N'\circ x_N^{-1}\bigr),
\]
that is, a product of chart transition maps of $\mathcal R$ and hence smooth. The lecturer leaves this as an exercise: ``it is immediately clear, and because it is immediately clear you may do it in an exercise: the chart transition maps are smooth, and from this follows that the configuration space is a smooth manifold''; the product of the topologies is mentioned as given \ts{15}{1}{25:14}--\ts{15}{1}{27:49}. \qedhere
\end{proof}

``Why did we define an \emph{atlas}? We have defined only one chart.'' Because the $x_j$ run through \emph{all} the charts we have on $\mathcal R$; with the single $x_j$'s going through all charts of $\mathcal R$ we obtain all the charts, hence the atlas \ts{15}{1}{26:00}--\ts{15}{1}{26:40}. ``Are there questions about the configuration space? Did you do this on the problem sheet? Come to me in the break'' \ts{15}{1}{27:49}.

% ==================================================================
\section{The motion of $N$ particles as one curve in $\mathcal Q$}\label{sec:mp-curve}

The second chapter has a handy title: describing the trajectories of $N$ particles in $\mathcal R$ as \emph{one curve in the configuration space $\mathcal Q$} \ts{15}{1}{28:14}--\ts{15}{1}{29:57}. Our starting point are $N$ trajectories $\gamma_1,\dots,\gamma_N\colon\R\to\mathcal R$, corresponding to the $N$ particles \ts{15}{1}{29:57}.

\begin{example}[Two particles on a line]\label{ex:mp-line}
Take a one-dimensional absolute space and on it two particles: one that flies a bit faster (it covers longer distances in the same times), and one which is slower: these are the trajectories $\gamma_1$ and $\gamma_2$ \ts{15}{1}{30:23}--\ts{15}{1}{31:29}.
\end{example}

We want to describe these curves as \emph{one} curve $\Gamma$ in the configuration space, $\Gamma\colon\R\to\mathcal Q$, ``from the real numbers to the configuration space'' \ts{15}{1}{31:55}--\ts{15}{1}{32:42}:
\begin{equation}\label{eq:mp-Gamma}
  \Gamma(\lambda)\;:=\;\bigl(\gamma_1(\lambda),\dots,\gamma_N(\lambda)\bigr)\ \in\mathcal Q ,
\end{equation}
``so we have indeed an $N$-tuple, and it is in $\mathcal Q$ because the individual curves each map into the absolute space'' \ts{15}{1}{33:02}--\ts{15}{1}{33:44}.

In the picture for the example: draw $\mathcal Q=\mathbb R\times\mathbb R$ as a plane, the absolute space of particle 1 on one axis and that of particle 2 on the other; the curve $\Gamma$ is the one in this plane whose coordinates are $(\gamma_1(\lambda),\gamma_2(\lambda))$ \ts{15}{1}{34:20}--\ts{15}{1}{35:41}.

Conversely, from a curve $\Gamma$ in $\mathcal Q$ we can recover the individual curves, with the \textbf{projections} $\mathrm{pr}_j\colon\mathcal Q\to\mathcal R$, $j=1,\dots,N$,
\begin{equation}\label{eq:mp-proj}
  \mathrm{pr}_j(p_1,\dots,p_N)\;:=\;p_j,\qquad\gamma_j\;=\;\mathrm{pr}_j\circ\Gamma ,
\end{equation}
(``the name projection comes from the fact that it projects the $j$-th component out of the tuple'') \ts{15}{1}{35:41}--\ts{15}{1}{38:02}. (The lecturer writes $p_j$ as an index-$j$ map $P_j$; with $\Gamma$ as the capital of $\gamma$.) Any questions? ``All clear, wonderful'' \ts{15}{1}{38:02}.

% ==================================================================
\section{Composite systems without interaction}\label{sec:mp-composite}

Point three: \textbf{composite systems without interaction between the systems}; we build several particles together as we just learned and want to set up the Lagrangian of such a system. (``For those who sit in the back and can read: it says $N$ particles.'') \ts{15}{1}{38:02}--\ts{15}{1}{39:43}.

How is such a composite system without interaction described? By a Lagrangian again, which is a map of the form $L\colon T\mathcal Q\to\R$: ``to the real numbers, that is the easy side, and from the tangent bundle of $\mathcal Q$'' \ts{15}{1}{39:43}--\ts{15}{1}{40:46}. What is the tangent bundle of $\mathcal Q$? We know what $T\mathcal R$ is, the tangent bundle over absolute space, but now over the configuration space:

\begin{theorem}[The tangent bundle of the configuration space; stated without proof]\label{thm:mp-TQ}
The tangent bundle of the configuration space is (in a canonical way) the Cartesian product of the tangent bundles,
\begin{equation}\label{eq:mp-TQ}
  T\mathcal Q\;\cong\;T\mathcal R\times\dots\times T\mathcal R\qquad(N\text{ factors}),
\end{equation}
so that a tangent vector to $\mathcal Q$ at $(p_1,\dots,p_N)$ is an $N$-tuple $(X_1,\dots,X_N)$ of vectors $X_j\in T_{p_j}\mathcal R$.
\end{theorem}

\begin{coursenote}[The proof is postponed (as said in class)]\label{cn:mp-proof}
``Whoever wants to see an exact proof should perhaps ask Frederik''; the lecturer tries to make it plausible with a picture: let $\mathcal R_1$ be a one-dimensional space and $\mathcal R_2$ ``a point'' (``because if we drew also a line the picture would become higher-dimensional, I already find three dimensions difficult on the blackboard''); the product is just the line, and at each point of that line sticks the tangent space of the point; but then the picture comes out ``doof'': ``forget it, this is silly, we leave the intuition away''. He promises that Frederik may do the proof in the next lecture, or at an office hour; ``it is silly like this, so let's leave it out'' \ts{15}{1}{41:31}--\ts{15}{1}{44:38}.
\end{coursenote}

(For orientation, my addition, not said in class: in the product chart $\omega$ of \eqref{eq:mp-omega}, the induced chart $\xi_\omega$ of $T\mathcal Q$ (\cref{def:lag-xi}) has the components $\bigl(x_1(p_1),\dots,x_N(p_N)\,;\,\mathrm dx_1(X_1),\dots,\mathrm dx_N(X_N)\bigr)$ in $\R^{2Nd}$; this is the chart of $T\mathcal R\times\dots\times T\mathcal R$ built from the $\xi_{x_j}$, which is what \eqref{eq:mp-TQ} expresses.)

So we have the map $L\colon T\mathcal Q\cong(T\mathcal R)^N\to\R$ and what we need is a rule for what it does. It should eat vector fields (here: vectors) over $\mathcal R$ and give back a real number; our idea (``habt ihr da vielleicht eine idee''): \emph{add up the individual one-particle Lagrangians} \ts{15}{1}{44:38}--\ts{15}{1}{45:43}.

\begin{definition}[New concepts: the Lagrangian of a composite system without interaction]\label{def:mp-Lsum}
Let $L_1,\dots,L_N\colon T\mathcal R\to\R$ be the Lagrangians of the $N$ individual particles (one-particle Lagrangians, each with its own mass and its own charges). The \textbf{Lagrangian of the composite system without interaction} is
\begin{equation}\label{eq:mp-Lsum}
  L\colon T\mathcal Q\to\R,\qquad L(X_1,\dots,X_N)\;:=\;L_1(X_1)+\dots+L_N(X_N)
\end{equation}
\ts{15}{1}{45:43}--\ts{15}{1}{47:13}.
\end{definition}
We insert into each one-particle Lagrangian its own vector: this is what we know since the last lecture, and it is a real number, so the sum is a real number: ``and with this we have also our first question settled: how do I construct from the Lagrangian of one particle the Lagrangian of $N$ particles'' \ts{15}{1}{46:05}--\ts{15}{1}{46:50}.

\begin{intuition}[Actions add (my remark, not said in class)]\label{int:mp-sum}
Since the lift of the curve $\Gamma$ in $T\mathcal Q$ is $\tilde\Gamma=(\tilde\gamma_1,\dots,\tilde\gamma_N)$, the action of the composite system is, by \eqref{eq:mp-Lsum}, $S[\Gamma]=\int\mathrm d\lambda\,\sum_jL_j(\tilde\gamma_j(\lambda))=\sum_jS_j[\gamma_j]$: the sum of the one-particle actions, as announced in \cref{int:lag-sum}.
\end{intuition}

\begin{example}[Two particles, two potentials]\label{ex:mp-two}
Look at two particles buzzing on the absolute space: the first with mass $m$, the second with mass $\tilde m$. There are two potentials $\Phi_1$ and $\Phi_2$ (of two conservative interaction fields), and it shall be that the first particle is coupled only to the first potential and not to the second: its charges are $c_1\neq0$, $c_2=0$; the second particle has charges $\tilde c_1=0$, $\tilde c_2\neq0$, so it couples only to the second \ts{15}{1}{47:43}--\ts{15}{1}{48:45}. Then the one-particle Lagrangians are
\begin{equation}\label{eq:mp-L12}
  \begin{aligned}
  L_1(X_1)&\;=\;\tfrac12\,m\,g_{\pi(X_1)}(X_1,X_1)-c_1\Phi_1\bigl(\pi(X_1)\bigr),\\
  L_2(X_2)&\;=\;\tfrac12\,\tilde m\,g_{\pi(X_2)}(X_2,X_2)-\tilde c_2\Phi_2\bigl(\pi(X_2)\bigr),
  \end{aligned}
\end{equation}
``again the metric at the foot point, the vector inserted twice, and next to it the potential'', and the Lagrangian of the whole system without interaction is the sum of the two:
\begin{equation}\label{eq:mp-L12sum}
  \begin{aligned}
  L(X_1,X_2)\;=\;&\tfrac12\,m\,g(X_1,X_1)+\tfrac12\,\tilde m\,g(X_2,X_2)\\&-c_1\Phi_1(\pi X_1)-\tilde c_2\Phi_2(\pi X_2),
  \end{aligned}
\end{equation}
each term evaluated at the corresponding foot point \ts{15}{1}{48:45}--\ts{15}{1}{51:52}. (The lecturer waits until everybody has copied; the middle blackboard was not yet written.)
\end{example}

% ==================================================================
\section{Interactions between the systems}\label{sec:mp-interaction}

Point four, the last item of material: a short chapter on \textbf{interactions}; the lecturer says that the next lecture will tell something more (perhaps with an example computed by Frederik, or one to be done on the problem sheet; ``we do not know exactly yet'') \ts{15}{1}{52:47}--\ts{15}{1}{53:29}.

\begin{definition}[New concepts: interaction between systems (Wechselwirkung zwischen Systemen)]\label{def:mp-interaction}
An \textbf{interaction} between the systems is a smooth map
\begin{equation}\label{eq:mp-I}
  I\colon T\mathcal Q\to\R ,
\end{equation}
that is, a Lagrangian again, \emph{without further restrictions} \ts{15}{1}{53:29}--\ts{15}{1}{54:46}.
\end{definition}
(The captions read ``von cq nach r'', which I read as $T\mathcal Q\to\R$ [?]. The name of the map in the lecture is a garbled ``and interaction''. How $I$ enters the Lagrangian is not said explicitly in this lecture: by analogy with \cref{int:lag-sum} one adds it, $L_{\text{total}}=L_1+\dots+L_N+I$ [my addition, not said in class].)

The point: the Lagrangians looked at so far, e.g.\ $L_1(X_1)$, depend only on the coordinates and velocities of \emph{one} particle, not of the other; now, without further restrictions, terms may also appear that depend on coordinates or velocities of \emph{two} (or more) particles. As an example, written already in coordinates (so the inverse chart map $\xi_\omega^{-1}$ has to be applied): a function $I$ of the arguments $q_1,q_2,\dot q_1,\dot q_2$, which could depend on all four, but we let it depend only on the two positions, say through a linear dependence on the distance (``irgendwie ein lineares Abstandsgesetz'', [?]) \ts{15}{1}{54:46}--\ts{15}{1}{56:22}.

\begin{coursenote}[Break (as said in class)]\label{cn:mp-break}
``Then I would say we meet again at five past eleven'': break; after it the promised worked example follows \ts{15}{1}{56:22}.
\end{coursenote}

% ==================================================================
\section{Example: two free particles}\label{sec:mp-example}

As the example take the simplest that we can think of: \emph{two particles} ($N=2$) which fly through the $d$-dimensional absolute space, without any external interaction field and without interacting with each other. ``This is the only thing that is difficult to imagine'' (the rest is easy) \ts{15}{1}{56:44}--\ts{15}{1}{58:16}. One particle has mass $m_1$ and the other mass $m_2$; they do not see each other, and see no interaction from outside \ts{15}{1}{58:16}.

We know that the Lagrangian is the sum of the one-particle Lagrangians:
\begin{equation}\label{eq:mp-Lfree}
  L(X_1,X_2)\;=\;\tfrac12\,m_1\,g_{\pi(X_1)}(X_1,X_1)\;+\;\tfrac12\,m_2\,g_{\pi(X_2)}(X_2,X_2).
\end{equation}
We want to apply the Euler--Lagrange equations, in the chart $\xi_\omega$ of $T\mathcal Q$ induced by the product chart $\omega$ of the configuration space (see \eqref{eq:mp-omega}), which has dimension $2d$ ($=Nd$ with $N=2$), so that $\xi_\omega$ goes into $\R^{4d}$ \ts{15}{1}{58:39}--\ts{15}{1}{59:42}. The Euler--Lagrange equations (\cref{thm:var-el}) read, for $\mu=1,\dots,2d$,
\begin{equation}\label{eq:mp-EL}
  \frac{\mathrm d}{\mathrm d\lambda}\Bigl(\partial_{2d+\mu}L_{(\omega)}\bigl(\xi_\omega\circ\tilde\Gamma(\lambda)\bigr)\Bigr)\;-\;\partial_\mu L_{(\omega)}\bigl(\xi_\omega\circ\tilde\Gamma(\lambda)\bigr)\;=\;0 ,
\end{equation}
where now it is a $2d$ in place of $d$ (``because first we must go through the $2d$ position coordinates of the particles, and the velocities, the dotted coordinates, are the ones after which we differentiate'') \ts{15}{1}{59:42}--\ts{15}{1}{1:00:48}. (The lecturer writes the chart representation as $L_{(\xi)}$ with a garbled symbol, $\xi_\omega$ here, and the lift of the curve as $\tilde\Gamma$.)

\subsection{The chart representation of $L$}\label{ssec:mp-chartrep}

For the Euler--Lagrange equations we need the chart representation $L_{(\omega)}=L\circ\xi_\omega^{-1}\colon\R^{4d}\to\R$, hence for given numbers. Call the $4d$ arguments
\begin{equation}\label{eq:mp-args}
  \bigl(\alpha^1,\dots,\alpha^d,\ \tilde\alpha^1,\dots,\tilde\alpha^d,\ \beta^1,\dots,\beta^d,\ \tilde\beta^1,\dots,\tilde\beta^d\bigr)
\end{equation}
(the first $2d$ numbers are the positions of the first and of the second particle; the last $2d$ are their velocity components; the tilde distinguishes the second particle; shorthand: $\alpha$ and $\tilde\alpha$ stand for all $d$ components) \ts{15}{1}{1:00:48}--\ts{15}{1}{1:04:53}. These numbers have nothing to do with the curve (``first differentiate, then insert the curve''). Since $L$ is a sum, so is $L_{(\omega)}$, and we use the structure of the product chart: the first particle is described by the chart $x_1$, the second by $x_2$, hence
\begin{equation}\label{eq:mp-Lomega}
  L_{(\omega)}(\alpha,\tilde\alpha,\beta,\tilde\beta)\;=\;L_{1,(x_1)}(\alpha,\beta)+L_{2,(x_2)}(\tilde\alpha,\tilde\beta),
\end{equation}
by the distributive law: ``this plus distributes over the circle'' \ts{15}{1}{1:04:53}--\ts{15}{1}{1:06:19}; the chart representation $\xi_\omega^{-1}$ ``is built from all the chart maps which we have on the absolute space'' and the term $L_1$ is interested only in the chart $x_1$, $L_2$ only in $x_2$ (``theoretically one could add the zeros''). The objects are the chart representations of the one-particle Lagrangians (as at the end of the last lecture); explicitly, by the definition of $L$,
\begin{equation}\label{eq:mp-Lchart}
  L_{1,(x_1)}(\alpha,\beta)=\tfrac12m_1\,g_{\rho\sigma}\!\bigl(x_1^{-1}(\alpha)\bigr)\beta^\rho\beta^\sigma,\qquad
  L_{2,(x_2)}(\tilde\alpha,\tilde\beta)=\tfrac12m_2\,g_{\rho\sigma}\!\bigl(x_2^{-1}(\tilde\alpha)\bigr)\tilde\beta^\rho\tilde\beta^\sigma ,
\end{equation}
where $g_{\rho\sigma}\circ x_j^{-1}$ denote the components of the metric in the chart $x_j$ (the lecturer notices an $x_2^{-1}$ that was forgotten and corrects it) \ts{15}{1}{1:06:19}--\ts{15}{1}{1:10:17}.

\subsection{The curve inserted; the shorthand}\label{ssec:mp-curve}

``Then we compute the Euler--Lagrange equations of all of this.'' We need the pieces of \eqref{eq:mp-EL}; we compute the two terms separately, term 1 and term 2, and one inserts successively all the definitions \ts{15}{1}{1:10:17}--\ts{15}{1}{1:11:25}.

\emph{The lift in the chart.} We must evaluate at $\xi_\omega\circ\tilde\Gamma(\lambda)$. The first $2d$ entries are the chart components of $\Gamma(\lambda)=(\gamma_1(\lambda),\gamma_2(\lambda))$, namely by the definition \eqref{eq:mp-omegacomp} of $\omega$ the numbers $x_1^m(\gamma_1(\lambda))$, then $x_2^m(\gamma_2(\lambda))$; the second $2d$ entries are the components of the tangent vector of the lifted curve, the $\lambda$-derivatives of these numbers. With the abbreviation, ``as always when we see such an $x$ after $\gamma$'',
\begin{equation}\label{eq:mp-q}
  q_1^m(\lambda)\;:=\;x_1^m\bigl(\gamma_1(\lambda)\bigr),\qquad q_2^m(\lambda)\;:=\;x_2^m\bigl(\gamma_2(\lambda)\bigr),
\end{equation}
this is
\begin{equation}\label{eq:mp-lift}
  \xi_\omega\circ\tilde\Gamma(\lambda)\;=\;\bigl(q_1(\lambda),\,q_2(\lambda),\,\dot q_1(\lambda),\,\dot q_2(\lambda)\bigr)
\end{equation}
(where $q_1(\lambda)$ stands for the $d$ numbers $q_1^1(\lambda),\dots,q_1^d(\lambda)$, and so on) \ts{15}{1}{1:11:25}--\ts{15}{1}{1:21:06}. The lecturer unwinds this step by step: insert the definition of $\xi_\omega$ (the entries $x_j^m(\pi(\cdot))$ and the components $\mathrm dx_j^m$), then the definition of $\tilde\Gamma$ as the lift of the curve $\Gamma=(\gamma_1,\gamma_2)$, then of $\omega$, giving the entries $x_1^m(\gamma_1(\lambda))$, \dots, $x_2^m(\gamma_2(\lambda))$ and their derivatives; the lecturer remarks that this is just the usual hyper-detailed bookkeeping, ``as one should have seen the last time'', and checks that everybody follows \ts{15}{1}{1:13:11}--\ts{15}{1}{1:21:42}.

\subsection{Derivatives with respect to the velocities}\label{ssec:mp-dbeta}

First \emph{term 1}, $\partial_{2d+\mu}L_{(\omega)}$ evaluated along the lift. Insert \eqref{eq:mp-Lomega} and \eqref{eq:mp-Lchart}; ``first differentiate, then insert the curve''. In the lecturer's words: the $\mu$-th derivative with respect to the velocity entries acts on $\beta^\rho\beta^\sigma$, which contains the $\beta$ twice, so one needs a product rule: it produces two terms (a $\delta^\rho_\mu$ in one and a $\delta^\sigma_\mu$ in the other), and since the metric is symmetric and the summation indices can be renamed, the two terms are equal: the factor $\tfrac12$ cancels \ts{15}{1}{1:21:42}--\ts{15}{1}{1:27:35}. For the derivative with respect to the velocity of the first particle, $\mu=1,\dots,d$ (the second particle's part does not depend on it):
\begin{equation}\label{eq:mp-p1}
  \partial_{2d+\mu}L_{(\omega)}\bigl(\xi_\omega\circ\tilde\Gamma\bigr)\;=\;m_1\,g_{\mu\sigma}\bigl(x_1^{-1}(q_1)\bigr)\,\dot q_1^{\sigma},\qquad\mu=1,\dots,d,
\end{equation}
and likewise, for the entries belonging to the second particle,
\begin{equation}\label{eq:mp-p2}
  \partial_{2d+d+\mu}L_{(\omega)}\bigl(\xi_\omega\circ\tilde\Gamma\bigr)\;=\;m_2\,g_{\mu\sigma}\bigl(x_2^{-1}(q_2)\bigr)\,\dot q_2^{\sigma},\qquad\mu=1,\dots,d,
\end{equation}
[the lecturer works out the first and says the second is entirely analogous \ts{15}{1}{1:28:00}--\ts{15}{1}{1:30:08}; the labelling of the entries by $\mu$ for the second particle, $\mu\to d+\mu$, is not written out in the lecture, [my addition]].

\begin{coursenote}[A student objects (as said in class)]\label{cn:mp-student}
While the lecturer carries out the derivative \eqref{eq:mp-p1}, a student objects to something (the captions are not understandable; probably: the product rule and the case distinction by which particle the index $\mu$ belongs to, so that the term of the other particle drops out [?]). The lecturer: ``yes, I know what you mean; I will look at it again, and Frederik will say something about it next week; I put a star here and I just keep computing, because something nice comes out, which is of course never a justification; but you are right, we have to make a small case distinction; in principle the right thing comes out'' \ts{15}{1}{1:36:50}--\ts{15}{1}{1:37:53}.
\end{coursenote}

\subsection{The $\lambda$-derivative}\label{ssec:mp-ddl}

Now \eqref{eq:mp-p1} must be differentiated with respect to $\lambda$ (``the first expression was only the thing inside the parenthesis'') \ts{15}{1}{1:37:53}--\ts{15}{1}{1:38:40}. Product rule with two terms: (i) $\lambda$ hits the second factor, which gives $\ddot q_1^\sigma$ (a second dot, ``we get a second point'' \ts{15}{1}{1:39:22}); (ii) $\lambda$ hits the metric components, which depend on $\lambda$ through $q_1(\lambda)$: the multidimensional chain rule (``which we derived on the first problem sheet'') gives $\partial_\nu g_{\mu\sigma}\,\dot q_1^\nu$, summed over $\nu$ ($\partial_\nu g$ meaning the derivative with respect to the $\nu$-th chart coordinate) \ts{15}{1}{1:38:40}--\ts{15}{1}{1:41:39}:
\begin{equation}\label{eq:mp-dterm1}
  \frac{\mathrm d}{\mathrm d\lambda}\,\partial_{2d+\mu}L_{(\omega)}\;=\;m_1\,g_{\mu\sigma}\,\ddot q_1^\sigma\;+\;m_1\,\frac{\partial g_{\mu\sigma}}{\partial x^\nu}\,\dot q_1^\nu\,\dot q_1^\sigma .
\end{equation}
\ts{15}{1}{1:41:39}--\ts{15}{1}{1:43:01}.

\subsection{Derivatives with respect to the positions}\label{ssec:mp-dalpha}

Term 2: $\partial_\mu L_{(\omega)}$ for $\mu=1,\dots,d$, the derivative with respect to the positions of the first particle; it only touches the metric components (``where do I have to apply the inverse chart map? $x_1\circ\gamma_1\ \dots$ the inverse chart gives back $\gamma_1$''):
\begin{equation}\label{eq:mp-dterm2}
  \partial_\mu L_{(\omega)}\bigl(\xi_\omega\circ\tilde\Gamma\bigr)\;=\;\tfrac12\,m_1\,\frac{\partial g_{\rho\sigma}}{\partial x^\mu}\,\dot q_1^\rho\dot q_1^\sigma ,
\end{equation}
all quantities evaluated at $x_1^{-1}(q_1(\lambda))=\gamma_1(\lambda)$ \ts{15}{1}{1:28:20}--\ts{15}{1}{1:29:25}, \ts{15}{1}{1:42:40}.

\subsection{Assembling the equations}\label{ssec:mp-assemble}

Put together, the Euler--Lagrange equation \eqref{eq:mp-EL} for $\mu=1,\dots,d$ becomes
\begin{equation}\label{eq:mp-elfree1}
  m_1\Bigl[\,g_{\mu\sigma}\,\ddot q_1^\sigma+\frac{\partial g_{\mu\sigma}}{\partial x^\nu}\dot q_1^\nu\dot q_1^\sigma-\frac12\frac{\partial g_{\rho\sigma}}{\partial x^\mu}\dot q_1^\rho\dot q_1^\sigma\Bigr]\;=\;0 ,
\end{equation}
where, as in \cref{ssec:var-verify}, we write the middle term twice with a factor $\tfrac12$ and exchange the two summation indices of one of them (using the symmetry in $\rho\leftrightarrow\sigma$, which is allowed since the product $\dot q^\rho\dot q^\sigma$ is symmetric): ``this is the trick that should be well known'' \ts{15}{1}{1:42:40}--\ts{15}{1}{1:45:11}:
\begin{equation}\label{eq:mp-elfree1b}
  m_1\Bigl[\,g_{\mu\sigma}\,\ddot q_1^\sigma+\frac12\Bigl(\frac{\partial g_{\mu\rho}}{\partial x^\sigma}+\frac{\partial g_{\mu\sigma}}{\partial x^\rho}-\frac{\partial g_{\rho\sigma}}{\partial x^\mu}\Bigr)\dot q_1^\rho\dot q_1^\sigma\Bigr]\;=\;0 .
\end{equation}
The same for the second particle: ``analogously every $1$ is replaced by a $2$'', and one gets the same equation with $m_2,q_2$, for the entries $d+1,\dots,2d$ \ts{15}{1}{1:44:48}.

Then the lecturer writes the result without the decorations, as it is found in the books \ts{15}{1}{1:45:11}: for each particle, $m_j$ times (the metric component times the acceleration plus the derivative terms in the dotted velocities) equal to zero, where he introduces for the bracket the symbol $g_{j\,\mu\rho,\sigma}$ meaning ``exactly what is above it'' (the notation of ``derivative with a comma'' from the problem sheet) \ts{15}{1}{1:45:32}--\ts{15}{1}{1:46:39}.

\subsection{Back to the free particles of experimental physics 1}\label{ssec:mp-back}

``You look at it and think: what is this good for? What did we start with? We started with two free particles which interact neither with each other nor with external fields; and now I stand here with this. How would two free particles look in the formalism of \emph{Experimentalphysik} 1?'' Can we conjure the two free particles of the introductory lecture out of this? \ts{15}{1}{1:46:39}--\ts{15}{1}{1:47:46}.

What does one do in experimental physics? One looks only at flat spaces. So we go to a \emph{flat} absolute space and, in addition, to a chart in which the metric takes the special form $g_{\rho\sigma}=\delta_{\rho\sigma}$ (``you probably know which form''); then all the derivatives of the metric vanish, since the metric is constant, and the equations \eqref{eq:mp-elfree1b} reduce to
\begin{equation}\label{eq:mp-newton}
  m_1\,\ddot q_1^{\,\mu}\;=\;0,\qquad m_2\,\ddot q_2^{\,\mu}\;=\;0,\qquad\mu=1,\dots,d\ (=1,2,3),
\end{equation}
for each particle separately. This is precisely what one learns in experimental physics 1 for \emph{two free particles without forces}: ``the same result as in the EP lecture'' \ts{15}{1}{1:47:46}--\ts{15}{1}{1:51:01}. (The lecturer's wording, in a flat Cartesian chart: $m_1\ddot q_1=0$, $m_2\ddot q_2=0$ for $\mu=1,2,3$.)

\begin{coursenote}[What the example shows, and the end of the lecture (as said in class and my remark)]\label{cn:mp-end}
The lecturer notes that with the lecture hour over he ``still has one minute but generously gives it to you'' \ts{15}{1}{1:51:22}. In terms of structure the example shows (my remark, not said in class) that for a system \emph{without interaction} the Euler--Lagrange equations of the sum decouple: the equations for the entries of particle $j$ contain only $q_j$, which is the same as the stationarity of each $S_j$ separately; coupling between the particles can come only from an interaction term $I$ of \cref{def:mp-interaction}.
\end{coursenote}

% ==================================================================
\section*{Key concepts}
\begin{keyconcepts}
\textbf{Retrieval questions.}
\begin{enumerate}
  \item Define the configuration space of $N$ particles. What is its dimension, and how does one build its smooth atlas from the atlas of absolute space?
  \item Write the product chart $\omega$ and its components; why are the chart transition maps smooth?
  \item How is the motion of $N$ particles described as one curve in $\mathcal Q$? How does one recover the individual trajectories?
  \item What is the tangent bundle of $\mathcal Q$? Write the Lagrangian of a composite system without interaction.
  \item What is the Lagrangian of two particles which couple to different potentials? Which terms contain the charges?
  \item What is an interaction between systems, and what may it depend on that a one-particle term may not?
  \item Compute the Euler--Lagrange equations for two free particles in a chart; what do they reduce to in a flat Cartesian chart?
\end{enumerate}
\textbf{Concepts.}
\begin{itemize}
  \item Configuration space $\mathcal Q=\mathcal R^N$ and its atlas (\cref{def:mp-config,def:mp-omega,thm:mp-smooth}, \eqref{eq:mp-config}, \eqref{eq:mp-omega}, \eqref{eq:mp-omegacomp}).
  \item The curve $\Gamma$ in $\mathcal Q$ and the projections (\eqref{eq:mp-Gamma}, \eqref{eq:mp-proj}).
  \item $T\mathcal Q\cong(T\mathcal R)^N$ and the sum Lagrangian (\cref{thm:mp-TQ,def:mp-Lsum}, \eqref{eq:mp-TQ}, \eqref{eq:mp-Lsum}).
  \item Interaction term (\cref{def:mp-interaction}, \eqref{eq:mp-I}).
  \item Two free particles: Euler--Lagrange equations in the product chart (\eqref{eq:mp-EL}, \eqref{eq:mp-elfree1b}, \eqref{eq:mp-newton}).
\end{itemize}
\end{keyconcepts}

\flushside
\end{paracol}
